\chapter{مقدمة}

\section{مقدمة عن مسائل القيم الابتدائية}
تُعرَّف مسائل القيم الابتدائية (Initial Value Problems) على أنها مسائل رياضية تتكون من معادلة تفاضلية عادية من الرتبة الأولى أو أعلى، مقترنة بشرط ابتدائي يحدد قيمة الحل عند نقطة معينة من مجال المتغير المستقل. الصيغة العامة لمسألة قيم ابتدائية من الرتبة الأولى تُعطى على الشكل
\[
\frac{dy}{dx} = f(x,y),
\]
مع شرط ابتدائي من الصورة
\[
y(x_0) = y_0.
\]

يمثل الشرط الابتدائي قيدًا أساسيًا يسمح بتحديد حل واحد من بين مجموعة لا نهائية من الحلول الممكنة للمعادلة التفاضلية. من الناحية النظرية، تُعد مبرهنة الوجود والوحدانية حجر الأساس في دراسة هذه المسائل، إذ تنص على وجود حل وحيد للمسألة ضمن مجال معين، شريطة تحقق شروط مناسبة على الدالة \(f(x,y)\)، مثل الاستمرارية وتحقيق شرط ليبشيتز.

تلعب مسائل القيم الابتدائية دورًا محوريًا في النمذجة الرياضية للظواهر الديناميكية، حيث تمثل الحالة الابتدائية للنظام نقطة الانطلاق لتطور الحل مع الزمن أو مع المتغير المستقل. ومن هنا تنبع أهميتها في مجالات متعددة مثل الميكانيكا، وانتقال الحرارة، والدوائر الكهربائية، والنماذج البيولوجية والاقتصادية.

\section{طرق حل مسائل القيم الابتدائية}
تختلف طرق حل مسائل القيم الابتدائية باختلاف طبيعة المعادلة التفاضلية ومدى تعقيدها. وبصورة عامة، يمكن تقسيم هذه الطرق إلى فئتين رئيسيتين: الطرق التحليلية والطرق العددية. يهدف هذا التقسيم إلى التمييز بين الطرق التي تعطي حلولًا دقيقة مغلقة، وتلك التي توفر حلولًا تقريبية قابلة للتطبيق العملي.

يعتمد اختيار طريقة الحل المناسبة على عدة عوامل، من أهمها قابلية المعادلة للحل التحليلي، ومتطلبات الدقة، والكلفة الحسابية، إضافة إلى طبيعة التطبيق العملي المرتبط بالمسألة.

\subsection{الطرق التحليلية}
تهدف الطرق التحليلية إلى إيجاد حل دقيق لمسألة القيم الابتدائية على شكل تعبير رياضي صريح أو ضمني. تعتمد هذه الطرق على أدوات التحليل الرياضي الكلاسيكي، مثل فصل المتغيرات، والمعادلات الخطية، وتحويلات لابلاس، وغيرها.

على سبيل المثال، تُعد المعادلة التفاضلية
\[
\frac{dy}{dx} = xy
\]
نموذجًا لمعادلة قابلة للفصل، ويمكن إرفاقها بشرط ابتدائي من الشكل
\[
y(0) = 1.
\]
يُستخدم هذا النوع من الأمثلة لإبراز مفهوم الحل التحليلي ودور الشرط الابتدائي في تحديد الحل الفريد، دون الخوض في تفاصيل الحساب.

وعلى الرغم من أن الطرق التحليلية توفر حلولًا دقيقة، إلا أن استخدامها يقتصر على فئات محدودة من المعادلات التفاضلية، إذ تفشل في التعامل مع معظم المعادلات غير الخطية أو المعقدة التي تظهر في التطبيقات الواقعية.

\subsection{الطرق العددية}
عندما يتعذر إيجاد حل تحليلي لمسألة القيم الابتدائية، تُستخدم الطرق العددية للحصول على تقريب عددي للحل عند نقاط منفصلة من مجال المتغير المستقل. تقوم هذه الطرق على فكرة استبدال المشتقة بتقريبات عددية، مما يؤدي إلى تحويل المعادلة التفاضلية إلى علاقات جبرية قابلة للحساب.

كمثال توضيحي، تُعد المسألة
\[
\frac{dy}{dx} = x + y, \quad y(0) = 1
\]
نموذجًا شائعًا لشرح فكرة الحل العددي، حيث يتم بناء متتالية من القيم التقريبية للحل باستخدام خوارزميات عددية مختلفة، مثل طريقة أويلر أو طرق رونج-كوتا.

تتميز الطرق العددية بمرونتها وقدرتها على التعامل مع نماذج رياضية معقدة، إلا أن دراسة خصائصها مثل الدقة والاستقرار والتقارب تُعد أمرًا ضروريًا لضمان موثوقية النتائج.

\section{تصنيف الطرق العددية}
يمكن تصنيف الطرق العددية المستخدمة في حل مسائل القيم الابتدائية وفقًا لعدد القيم السابقة التي تعتمد عليها في حساب القيمة الجديدة للحل. ووفقًا لهذا المعيار، تنقسم الطرق العددية إلى طرق وحيدة الخطوة وطرق متعددة الخطوات.

يُعد هذا التصنيف أساسيًا في التحليل النظري للطرق العددية، إذ يؤثر بشكل مباشر على خصائص الاستقرار والكفاءة الحسابية لكل طريقة.

\subsection{طرق وحيدة الخطوة}
تعتمد طرق وحيدة الخطوة على قيمة الحل عند نقطة واحدة فقط من أجل حساب القيمة التقريبية في النقطة التالية. الصيغة العامة لهذه الطرق يمكن تمثيلها على الشكل
\[
y_{n+1} = y_n + h\,\Phi(x_n,y_n,h),
\]
حيث \(h\) تمثل طول الخطوة، و\(\Phi\) دالة تحدد آلية التقريب المستخدمة.

تُعد هذه الطرق سهلة التطبيق وبسيطة من الناحية المفاهيمية، كما أنها لا تتطلب سوى قيمة ابتدائية واحدة لبدء الحساب. إلا أن تحقيق دقة عالية باستخدام هذه الطرق قد يتطلب اختيار خطوة صغيرة نسبيًا، مما يزيد من الكلفة الحسابية.

\subsection{طرق متعددة الخطوات}
تعتمد طرق متعددة الخطوات على أكثر من قيمة سابقة للحل من أجل حساب القيمة الجديدة. الصيغة العامة للطرق الخطية متعددة الخطوات تُعطى على الشكل
\[
\sum_{j=0}^{k} \alpha_j \, y_{n+j}
=
h \sum_{j=0}^{k} \beta_j \, f(x_{n+j}, y_{n+j}),
\]
حيث تمثل \(\alpha_j\) و\(\beta_j\) معاملات ثابتة تحدد نوع الطريقة، و\(k\) عدد الخطوات المستخدمة.

تعكس هذه الصيغة العامة الفكرة الأساسية للطرق متعددة الخطوات، والمتمثلة في الاستفادة من معلومات الحل عند عدة نقاط سابقة لتحسين الدقة وتقليل عدد العمليات الحسابية. وتُعد طرق آدامز-باشفورث وآدامز-مولتون أمثلة شهيرة على هذا النوع من الطرق.

تمتاز الطرق متعددة الخطوات بكفاءتها العالية، إلا أنها تتطلب قيمًا ابتدائية متعددة لبدء العملية الحسابية، والتي غالبًا ما يتم الحصول عليها باستخدام طرق وحيدة الخطوة. كما أن تحليل الاستقرار الصفري والتقارب يُعد عنصرًا جوهريًا في دراسة هذه الطرق.
